\chapter[Richiami: LP, QP, NLP e analisi di sensitività]{Richiami: programmazione lineare,
quadratica e non lineare}
\label{cap:richiami}

Questo capitolo raccoglie gli strumenti teorici usati in tutta la dispensa: la dualità
lineare, la convessità, le condizioni KKT e il protocollo di analisi di sensitività.
Chi li conosce già può passare direttamente al capitolo~\ref{cap:pythongurobi}.

\section{Programmazione lineare e dualità}

\begin{modello}[Coppia primale-duale]
\vspace{-1ex}
\begin{align*}
\text{(P)}\quad
& \min\; \sum_{j=1}^{n} c_j\, x_j
&& \text{s.t.}\;\; \sum_{j=1}^{n} a_{kj}\, x_j \ge b_k \;\; \forall k = 1, \dots, m;
   \quad x_j \ge 0 \;\; \forall j \\
\text{(D)}\quad
& \max\; \sum_{k=1}^{m} b_k\, y_k
&& \text{s.t.}\;\; \sum_{k=1}^{m} a_{kj}\, y_k \le c_j \;\; \forall j = 1, \dots, n;
   \quad y_k \ge 0 \;\; \forall k
\end{align*}
\end{modello}

I \emph{dati} sono i costi $c_j \in \R$, i coefficienti $a_{kj} \in \R$ e i termini
noti $b_k \in \R$; le \emph{variabili} sono le $x_j$ del primale e, nel duale, una
$y_k$ per ciascuno degli $m$ vincoli del primale. (In forma compatta:
$\min \vet c\T\vet x$ s.t. $\mat A \vet x \ge \vet b$, $\vet x \ge \vet 0$ --- la
useremo raramente.)
All'ottimo i due valori coincidono (\emph{dualità forte}) e la variabile duale $y_k$
misura di quanto cambia il valore ottimo se il termine noto $b_k$ aumenta di una unità:
è il \textbf{prezzo ombra} della risorsa $k$.

\begin{esempio}[Dualità in un problema di produzione $2\times 2$]
Un'azienda produce due prodotti con margini $30$ e $50$ \euro/unità:
\begin{align*}
\max\;& 30x_A + 50x_B \\
\text{s.t.}\;\; & x_A + 3x_B \le 90 \quad\text{(ore di lavorazione)}\\
& 2x_A + x_B \le 80 \quad\text{(kg di materia prima)}\\
& x_A, x_B \ge 0
\end{align*}
\textbf{Passo 1 — soluzione grafica/algebrica.} All'ottimo entrambi i vincoli sono attivi
(le rette si intersecano in un vertice). Risolviamo il sistema:
\begin{gather*}
\begin{cases} x_A + 3x_B = 90 \\ 2x_A + x_B = 80 \end{cases}
\;\Rightarrow\;
x_A = 90 - 3x_B,\quad 2(90 - 3x_B) + x_B = 80 \\
\Rightarrow\; 180 - 5x_B = 80
\;\Rightarrow\; x_B = 20,\; x_A = 30.
\end{gather*}
Valore ottimo: $z^* = 30\cdot 30 + 50 \cdot 20 = 1900$~\euro.

\textbf{Passo 2 — prezzi ombra.} I duali $(y_1, y_2)$ soddisfano le condizioni
$\mat A\T \vet y = \vet c$ sulle variabili in base:
\[
\begin{cases} y_1 + 2y_2 = 30 \\ 3y_1 + y_2 = 50 \end{cases}
\;\Rightarrow\; y_1 = 30 - 2y_2,\;\; 3(30 - 2y_2) + y_2 = 50 \;\Rightarrow\;
y_2 = 8,\; y_1 = 14.
\]
Un'ora di lavorazione in più vale \textbf{14~\euro}, un kg di materia prima \textbf{8~\euro}.

\textbf{Passo 3 — verifica della dualità forte.}
$\vet b\T \vet y = 90 \cdot 14 + 80 \cdot 8 = 1260 + 640 = 1900 = z^*$. \hfill$\checkmark$
\end{esempio}

\begin{spiegazione}
Due letture da ricordare per il laboratorio:
\begin{itemize}
  \item \textbf{Prezzo ombra} ($y_k$, attributo \texttt{Pi} in Gurobi): valore marginale
        della risorsa $k$. Vale solo finché la base ottima non cambia, cioè per
        $b_k$ nell'intervallo \texttt{SARHSLow}--\texttt{SARHSUp}. Un vincolo non attivo
        ha sempre prezzo ombra nullo.
  \item \textbf{Costo ridotto} (\texttt{RC}): per una variabile a zero, di quanto deve
        migliorare il suo coefficiente in obiettivo perché convenga attivarla.
\end{itemize}
Attenzione ai segni: Gurobi riporta i duali con la convenzione ``derivata del valore
ottimo rispetto al termine noto''. In un problema di \emph{minimo} un vincolo
$\le$ di capacità ha duale $\le 0$: allentare la capacità riduce il costo.
\end{spiegazione}

\section{Convessità e programmazione quadratica}

Un problema $\min f(\vet x)$ su un insieme ammissibile $X$ è \textbf{convesso} se $f$ è
convessa e $X$ è convesso. La proprietà che useremo di continuo:

\begin{insight}[Perché la convessità conta]
In un problema convesso ogni minimo locale è globale: il solver può \emph{certificare}
l'ottimalità. Nei problemi non convessi (ad esempio il pricing con termine bilineare
$p\cdot q$ del capitolo~\ref{cap:pricing}) un solver locale garantisce solo un ottimo
locale; Gurobi tratta i QP non convessi con tecniche globali, ma il costo computazionale
cresce.
\end{insight}

Un \textbf{QP} ha obiettivo $\tfrac12 \sum_{i}\sum_{j} q_{ij}\, x_i x_j + \sum_j c_j x_j$ (in forma compatta $\tfrac12 \vet x\T \mat Q\, \vet x + \vet c\T \vet x$) e vincoli lineari; è convesso se
e solo se $\mat Q \succeq 0$ (semidefinita positiva). Esempi nella dispensa: la varianza di
portafoglio $\sum_i\sum_j q_{ij} x_i x_j$ con $\mat Q = (q_{ij})$ matrice di covarianza (capitolo~\ref{cap:markowitz});
la norma $\|\vet w\|^2$ della SVM (capitolo~\ref{cap:svm}); le penalità di smoothing
$\sum_{t \in T,\, t > 1} (x_t - x_{t-1})^2$ del capitolo~\ref{cap:produzione}.

\section{Programmazione non lineare e condizioni KKT}

Le condizioni di ottimalità dei problemi non lineari vincolati prendono il nome di
\textbf{condizioni KKT} (da Karush, Kuhn e Tucker, che le formularono).
Per il problema $\min f(\vet x)$ s.t. $g_i(\vet x) \le 0$ per $i = 1, \dots, m$ e $h_j(\vet x) = 0$ per $j = 1, \dots, q$, la Lagrangiana è
\[
L(\vet x, \vet\lambda, \vet\nu) = f(\vet x) + \sum_{i=1}^{m} \lambda_i\, g_i(\vet x) + \sum_{j=1}^{q} \nu_j\, h_j(\vet x),
\qquad \lambda_i \ge 0 .
\]

\begin{modello}[Condizioni KKT]
In un punto di ottimo regolare esistono moltiplicatori $\vet\lambda \ge \vet 0$, $\vet\nu$ tali che:
\[
\nabla f(\vet x^*) + \textstyle\sum_{i=1}^{m} \lambda_i \nabla g_i(\vet x^*) + \sum_{j=1}^{q} \nu_j \nabla h_j(\vet x^*) = \vet 0,
\qquad
\lambda_i\, g_i(\vet x^*) = 0 \quad \forall i = 1, \dots, m \quad\text{(complementarietà)}.
\]
Se il problema è convesso, le KKT sono anche \emph{sufficienti}: un punto che le
soddisfa è un ottimo globale.
\end{modello}

\begin{esempio}[KKT in due variabili, svolto per intero]
$\min\; x^2 + y^2$ s.t. $x + y \ge 4$.

\textbf{Passo 1.} Riscriviamo il vincolo in forma standard: $g(x,y) = 4 - x - y \le 0$.

\textbf{Passo 2.} Stazionarietà della Lagrangiana $L = x^2 + y^2 + \lambda(4 - x - y)$:
\[
\frac{\partial L}{\partial x} = 2x - \lambda = 0, \qquad
\frac{\partial L}{\partial y} = 2y - \lambda = 0
\;\;\Rightarrow\;\; x = y = \lambda/2 .
\]

\textbf{Passo 3.} Complementarietà: se il vincolo fosse inattivo ($\lambda = 0$)
avremmo $x = y = 0$, che viola $x + y \ge 4$: impossibile. Quindi il vincolo è attivo:
$x + y = 4 \Rightarrow x = y = 2$, $\lambda = 4$.

\textbf{Passo 4.} Lettura del moltiplicatore: se il termine noto passa da $4$ a
$4 + \varepsilon$, il valore ottimo $f^* = 8$ cresce di circa
$\lambda\varepsilon = 4\varepsilon$. Verifica: $f^*(4+\varepsilon) =
2\left(\frac{4+\varepsilon}{2}\right)^2 = 8 + 4\varepsilon + \varepsilon^2/2$.
\hfill$\checkmark$
\end{esempio}

\begin{spiegazione}
I moltiplicatori KKT generalizzano i prezzi ombra dell'LP: misurano la variazione
marginale del valore ottimo rispetto al termine noto del vincolo. Nel laboratorio,
quando il solver non fornisce i duali (QP non convessi, solver locali), stimiamo il
moltiplicatore \emph{per perturbazione}: si aumenta il termine noto di
$\varepsilon$, si riottimizza e si calcola il rapporto incrementale. Questa verifica
numerica è essa stessa un esercizio ricorrente.
\end{spiegazione}

\section{L'analisi di sensitività negli LP e i prezzi ombra}
\label{sec:sensitivita-lp}

Risolvere un LP dà una soluzione ottima; l'\textbf{analisi di sensitività} risponde
alla domanda successiva, quella che interessa davvero un decisore: \emph{come cambia
l'ottimo se cambiano i dati?} Negli LP gran parte della risposta arriva gratis dal
solver, sotto forma di tre numeri per ogni vincolo e variabile:

\begin{itemize}
  \item il \textbf{prezzo ombra} di un vincolo (variabile duale, \texttt{Pi}):
        di quanto migliora il valore ottimo se il termine noto del vincolo aumenta di
        una unità --- ``quanto pagherei per un'unità in più di quella risorsa'';
  \item il \textbf{range di validità} del termine noto
        (\texttt{SARHSLow}--\texttt{SARHSUp}): l'intervallo in cui quel prezzo ombra
        resta esatto (fuori, la base ottima cambia e va rifatta l'ottimizzazione);
  \item il \textbf{costo ridotto} di una variabile a zero (\texttt{RC}): di quanto
        deve migliorare il suo coefficiente in obiettivo perché convenga attivarla.
\end{itemize}

\begin{esempio}[Prezzi ombra in pratica, sull'LP $2\times2$]
Riprendiamo l'esempio di produzione: ottimo $1900$~\euro{} con duali
$y_{\text{ore}} = 14$ e $y_{\text{mat}} = 8$.

\textbf{Lettura.} Un'ora di lavorazione in più vale $14$~\euro{} di ricavo aggiuntivo;
un kg di materia prima $8$~\euro. Se un fornitore offre ore di straordinario a
$12$~\euro/ora, conviene comprarle ($14 > 12$).

\textbf{Verifica per perturbazione.} Rirrisolvendo con $91$ ore: ricavo $1914$,
esattamente $+14$. Con $100$ ore: $2040 = 1900 + 10 \cdot 14$: il prezzo ombra è
ancora esatto perché $100$ è dentro il range di validità $[40, 240]$.

\textbf{Dove smette di valere.} A $240$ ore il vincolo delle ore smette di essere
la risorsa scarsa (la base cambia): oltre, ogni ora aggiuntiva vale meno di 14 e va
ricalcolato tutto. Ecco perché nel laboratorio ogni prezzo ombra si cita
\emph{insieme} al suo intervallo di validità.

\textbf{Vincoli non attivi.} Se le ore disponibili fossero $300$ (vincolo non
attivo all'ottimo), il loro prezzo ombra sarebbe $0$: una risorsa che avanza non
vale nulla al margine.
\end{esempio}

\section{Il protocollo di analisi di sensitività}
\label{sec:protocollo}

L'analisi di sensitività è parte del modello, non un'appendice. In ogni esercitazione
seguiamo lo stesso protocollo in sei passi:

\begin{center}\small
\begin{tabular}{lp{9.2cm}}
\toprule
1. Scenario base & risolvere, verificare l'ammissibilità, identificare i vincoli attivi. \\
2. One-at-a-time & variare un parametro chiave su una griglia; grafico valore/decisioni. \\
3. Prezzi ombra & confrontare la previsione marginale del duale con una ri-ottimizzazione
                  dopo una piccola perturbazione. \\
4. Scenari & costruire almeno tre scenari coerenti (pessimistico, centrale, ottimistico). \\
5. Trade-off & variare un peso o una soglia per tracciare una frontiera
               (costo-servizio, rischio-rendimento, costo-emissioni). \\
6. Stabilità & segnalare cambi bruschi del mix ottimo e parametri stimati fragili. \\
\bottomrule
\end{tabular}
\end{center}

\begin{esercizio}[dualità]
Nel problema $2\times2$ dell'esempio, il fornitore offre 10 ore di lavorazione extra a
12~\euro/ora. Conviene? Fino a che prezzo orario conviene? E fino a quante ore vale il
prezzo ombra di 14~\euro? (Suggerimento: ricalcolare il vertice ottimo quando il
vincolo delle ore si sposta a $90+\Delta$ e trovare il $\Delta$ oltre il quale
la base cambia.)
\end{esercizio}

\begin{esercizio}[KKT]
Risolvere con le KKT: $\min\; (x-3)^2 + (y-1)^2$ s.t. $x + 2y \le 2$, $x, y \ge 0$.
Dire quali vincoli sono attivi all'ottimo e interpretare i moltiplicatori.
\end{esercizio}

\begin{esercizio}[convessità]
Per quali valori del parametro $\gamma$ la funzione
$f(x_1, x_2) = x_1^2 + \gamma x_1 x_2 + x_2^2$ è convessa? (Scrivere $Q$ e imporre
$Q \succeq 0$.)
\end{esercizio}
